\newcommand{\auditoriaid}{A01 -- Portada, Resumen y Alcance}
% Preambulo comun de las auditorias del informe E1 (revision_e1/)
\documentclass[11pt,a4paper]{article}
\usepackage[utf8]{inputenc}
\usepackage[T1]{fontenc}
\usepackage[spanish,es-noshorthands]{babel}
\usepackage{lmodern}
\usepackage{microtype}
% Sin patrones de guionado en espanol en esta instalacion (ver A10): se
% desactiva el guionado automatico para evitar cortes con reglas del ingles.
\hyphenpenalty=10000
\exhyphenpenalty=10000
\usepackage[a4paper,margin=2.2cm]{geometry}
\usepackage{booktabs,array,longtable}
\usepackage{graphicx,caption,subcaption}
\usepackage{xcolor}
\usepackage{enumitem}
\usepackage{fancyhdr}
\usepackage{amsmath}
\usepackage{tikz}
\usetikzlibrary{arrows.meta,positioning}
\usepackage{natbib}
\usepackage{xurl}
\usepackage[hidelinks]{hyperref}
\urlstyle{tt}
\newcommand{\archivo}[1]{\url{#1}}
\newcolumntype{L}[1]{>{\raggedright\arraybackslash}p{#1}}
\definecolor{alta}{RGB}{180,30,30}
\definecolor{media}{RGB}{200,120,0}
\definecolor{baja}{RGB}{60,110,160}
\definecolor{cajafondo}{RGB}{245,247,250}
\newcommand{\sevA}{\textcolor{alta}{\textbf{Alta}}}
\newcommand{\sevM}{\textcolor{media}{\textbf{Media}}}
\newcommand{\sevB}{\textcolor{baja}{\textbf{Baja}}}
\setlength{\parskip}{0.35em}
\setlength{\emergencystretch}{3em}
\setlength{\parindent}{0pt}
\pagestyle{fancy}
\fancyhf{}
\fancyhead[L]{\small Auditoría del Informe del Entregable~1 -- \auditoriaid}
\fancyhead[R]{\small \thepage}
\renewcommand{\headrulewidth}{0.4pt}
% Entorno de hallazgos: ID | Severidad | Ubicacion | Hallazgo y evidencia | Correccion
\newenvironment{hallazgos}{%
\small
\begin{longtable}{@{}L{1.15cm}L{1.25cm}L{1.9cm}L{6.0cm}L{\dimexpr\textwidth-10.3cm-10.6\tabcolsep\relax}@{}}
\toprule
\textbf{ID} & \textbf{Sev.} & \textbf{Ubicación} & \textbf{Hallazgo y evidencia} & \textbf{Corrección aplicada} \\
\midrule\endfirsthead
\toprule
\textbf{ID} & \textbf{Sev.} & \textbf{Ubicación} & \textbf{Hallazgo y evidencia} & \textbf{Corrección aplicada} \\
\midrule\endhead
\bottomrule\endfoot
}{\end{longtable}}
% Macros compartidas con el informe revisado (las secciones propuestas las usan)
% Macros compartidas entre el informe revisado y las auditorias.
% Cifras verificadas contra los archivos del repositorio (ver A01/A04/A07):
\newcommand{\nUniverso}{102}   % source_id CMIP6 con tos/Omon en ESGF (00b, model_availability_priority.csv)
\newcommand{\nSeed}{41}        % publican los cuatro experimentos (config/models_seed_cmip6.csv)
\newcommand{\nSel}{40}         % completos + QC PASS (models_inventory_final.csv)
\newcommand{\nComb}{160}       % combinaciones modelo-experimento PASS (qc_report.csv)
\newcommand{\nNoEnc}{44}       % config/models_no_encontrado_44.csv
\newcommand{\nParcial}{18}     % 102 - 40 - 44
% Fecha de la portada: fija (no \today), editable en un solo lugar.
\newcommand{\fechaentrega}{28 de septiembre de 2026}
% Estado de codigo que documenta el E1 (antes de los cambios del E2)
\newcommand{\commitEone}{\texttt{8ee6184}}

% En las auditorias el texto propuesto se muestra fuera del informe: las
% referencias cruzadas se imprimen con el nombre de la seccion destino.
\makeatletter
\def\etq#1#2{\expandafter\def\csname etq@#1\endcsname{#2}}
\etq{sec:resumen}{Resumen}
\etq{sec:alcance}{Alcance}
\etq{sec:marco}{Revisión científica}
\etq{sec:antecedentes}{Antecedentes}
\etq{sec:referencias}{Literatura sobre TOE}
\etq{tab:referencias}{bibliografía TOE}
\etq{sec:similitudes}{Niño~3.4 frente a Niño~1+2}
\etq{sec:toe_metodos}{Elección de los métodos de TOE}
\etq{tab:metodos_toe}{métodos TOE}
\etq{eq:szabo_signal}{ecuación señal M1}
\etq{eq:szabo_V}{ecuación V(t)}
\etq{eq:szabo_T}{ecuación T(t)}
\etq{eq:szabo_toe}{ecuación TOE M1}
\etq{eq:tod}{ecuación ToD}
\etq{sec:datos}{Datos}
\etq{fig:cmip6_contribution}{contribución institucional}
\etq{sec:tabla_modelos}{Modelos seleccionados}
\etq{tab:modelos}{modelos}
\etq{sec:desviacion}{Selección de modelos}
\etq{sec:dominios_periodos}{Dominios y periodos}
\etq{fig:region_nino}{regiones Niño}
\etq{sec:arquitectura}{Arquitectura del pipeline}
\etq{fig:flujograma}{Flujograma}
\etq{sec:scripts}{Descripción de los pasos}
\etq{sec:paso00b}{Paso 00b}
\etq{sec:paso01}{Paso 01}
\etq{sec:paso02}{Paso 02}
\etq{sec:paso04}{Paso 04}
\etq{sec:paso05}{Paso 05}
\etq{sec:paso06}{Paso 06}
\etq{sec:paso07}{Paso 07}
\etq{sec:reporte}{Reporte de estado}
\etq{sec:resultados}{Resultados}
\etq{fig:qc}{matriz de estado}
\etq{tab:cuentas}{conteo de modelos}
\etq{fig:mapas2d}{mapas 2D}
\etq{fig:series}{series}
\etq{fig:boxplot}{boxplots}
\etq{sec:roadmap}{Próximos pasos}
\etq{tab:roadmap}{próximos pasos}
\etq{sec:conclusiones}{Conclusiones}
\etq{sec:recomendaciones}{Recomendaciones}
\makeatother

\makeatletter
\AtBeginDocument{\renewcommand{\ref}[1]{\@ifundefined{etq@#1}{\textit{[#1]}}{\textit{\csname etq@#1\endcsname}}}}
\makeatother
% Texto propuesto: secciones degradadas un nivel y en caja
\newenvironment{propuesto}{%
  \par\medskip\noindent\textcolor{baja}{\rule{\textwidth}{0.8pt}}\par
  \noindent{\small\textbf{Inicio del texto propuesto para el informe revisado}}\par\medskip
  \begingroup
  \let\section\subsection\let\subsection\subsubsection\let\subsubsection\paragraph
}{%
  \endgroup
  \par\noindent{\small\textbf{Fin del texto propuesto}}\par
  \noindent\textcolor{baja}{\rule{\textwidth}{0.8pt}}\par\medskip}
% Recuadro de actualizacion posterior (decisiones del autor)
\newcommand{\actualizacion}[1]{\par\medskip\noindent\fcolorbox{media}{media!8}{\parbox{\dimexpr\linewidth-2\fboxsep-2\fboxrule\relax}{\small\textbf{Actualización del 29-09-2026 (decisiones del autor).} #1}}\par\medskip}

\begin{document}

\begin{center}
{\Large\bfseries Auditoría A01: Portada, Resumen y Alcance}\\[0.4em]
{\normalsize Informe del Entregable~1 (\texttt{informe/informe\_e1\_smnv3.tex}, commit \texttt{330fd2e})}\\[0.2em]
{\small Revisión: \fechaentrega}
\end{center}

\section{Objeto}

Esta auditoría revisa la portada, el Resumen (Sección~1) y el Alcance (Sección~2) del informe del Entregable~1. Cada afirmación se contrastó contra tres fuentes: el texto de los términos de referencia (TdR, revisión \texttt{rev\_02092026}), los archivos de datos producidos por el \emph{pipeline} y el código en el estado que documenta el entregable (commit \texttt{8ee6184}, último anterior a los cambios introducidos para el Entregable~2). El resultado es una lista de hallazgos con su corrección y el texto ya corregido que pasa al informe revisado.

\actualizacion{El informe no cita commits. El Resumen ya no remite a secciones de pasos individuales: las seis correcciones se detallan en el anexo.}

\section{Fuentes contrastadas}

\begin{itemize}[leftmargin=1.2em]
\item TdR: numerales 5, 6 (actividades a--e), 8 y 10 (plazos de entregables).
\item \archivo{informe/model_availability_priority.csv}: 102 filas, todas con \texttt{tiene\_tos=1}; es el universo completo de \texttt{source\_id} que publican \texttt{tos}/\texttt{Omon} en el índice de ESGF.
\item \archivo{config/models_seed_cmip6.csv}: 41 modelos con los cuatro experimentos publicados.
\item \archivo{data/processed/models_inventory_final.csv} y \archivo{qc_report.csv}: 40 modelos \texttt{selected=True}; 160 filas, todas \texttt{PASS}.
\item \archivo{config/models_no_encontrado_44.csv}: 44 modelos.
\item Código de los pasos 00b, 01, 04, 05, 06 y 07 en el commit \texttt{8ee6184}; historial de \texttt{git} (commits \texttt{dc0bf9d}, \texttt{c07f4c5}, \texttt{79a9786}, \texttt{23cf218}).
\end{itemize}

\section{Verificación de cifras}

\begin{table}[h]
\centering\small
\begin{tabular}{@{}lrl@{}}
\toprule
Conjunto & N & Fuente \\
\midrule
Universo con \texttt{tos}/\texttt{Omon} en ESGF & 102 & \archivo{model_availability_priority.csv} \\
Publican los cuatro experimentos (\emph{seed}) & 41 & \archivo{models_seed_cmip6.csv} \\
Completos y aprobados en QC & 40 & \archivo{models_inventory_final.csv} \\
Combinaciones modelo-experimento PASS & 160 & \archivo{qc_report.csv} \\
Sin algún experimento en ninguna fuente & 44 & \archivo{models_no_encontrado_44.csv} \\
Disponibilidad o descarga incompleta & 18 & 17 fuera del \emph{seed} + 1 del \emph{seed} sin completar \\
\midrule
Control: $40 + 44 + 18$ & 102 & \\
\bottomrule
\end{tabular}
\end{table}

Las cifras 40/102 y 160 del informe son correctas. La descomposición del universo que da el Resumen (``46 de un barrido inicial del vocabulario CMIP6 completo más 56 citados en la literatura del proyecto'') no lo es: el paso~00b enumera directamente todos los \texttt{source\_id} con \texttt{tos}/\texttt{Omon} publicados en ESGF, y el propio código documenta que ese universo es de 102 modelos. El desglose 46\,+\,56 proviene de las versiones~1 y~2 del informe, cuando el punto de partida era una lista curada; sobrevivió a la reescritura de la versión~3.

\section{Hallazgos}

\begin{hallazgos}
A01-1 & \sevA & Resumen, 2.º párrafo & ``Del universo de 102 modelos candidatos (46 de un barrido inicial \ldots\ más 56 citados en la literatura)''. Falso en el estado actual: los 102 son el universo ESGF completo con \texttt{tos}/\texttt{Omon} (ver tabla). El desglose repite el diseño de v1/v2. & Se describe el universo como los 102 modelos que publican \texttt{tos}/\texttt{Omon} en ESGF, con los niveles 41 (\emph{seed}) y 40 (seleccionados). \\
A01-2 & \sevM & Resumen, 2.º párrafo & ``el resto quedó excluido por descarga incompleta o no disponible'' sin cifras, mientras que Resultados sí da 44 y deja el resto indeterminado (``un subconjunto adicional''). & Se dan las dos cifras, 44 y 18, que cierran el total. \\
A01-3 & \sevM & Resumen, 2.º párrafo & Enumera cinco problemas corregidos (FGOALS-f3-L y ``cuatro bugs de robustez'') en una sola oración de más de 170 palabras. Omite la corrección de unidades por valor (GISS-E2-1-G, Kelvin declarado como \texttt{degC}), que sí está en el código (\texttt{fix\_units} del paso~04) y en la Sección del paso~04. & Párrafo propio con los seis problemas, una línea cada uno y una referencia a su sección. \\
A01-4 & \sevA & Resumen, 3.er párrafo & ``desarrolla el marco teórico exigido por los TdR: la partición de \ldots\ Hawkins, Szabó, Bruno, Álvarez, Takahashi''. El TdR (numerales 5 y 6) no nombra autores; exige la revisión de metodologías de TOE y, para el Entregable~4, la partición de incertidumbre, la relación señal-ruido y el TOE. & ``El informe desarrolla la revisión científica: \ldots''. La elección de autores queda como decisión técnica del equipo. \\
A01-5 & \sevM & Alcance, 1.er párrafo & El sujeto ``Los Términos de referencia'' concuerda con ``define'' en singular; ``Las tareas (i) son'' designa una sola tarea. & ``definen''; ``La tarea~(i) es''. \\
A01-6 & \sevM & Alcance & El informe afirma reproducibilidad total, pero después del E1 el repositorio cambió la ventana de descarga (\texttt{79a9786}: de 100°E--70°O/20°S--20°N a global/30°S--30°N) y el techo del QC (\texttt{23cf218}: 39 → 45\,°C). Desde \texttt{main} no se reconstruye el estado del informe. & Se agregó una frase con el commit del E1. \textit{Actualización 29-09: por decisión del autor el informe no menciona commits; la frase se retiró y los datos descritos son los de la ventana actual.} \\
A01-7 & \sevM & Alcance & No se aclara que el QC de este entregable verifica integridad y homogeneidad, no desempeño; la aclaración aparecía recién en Recomendaciones. & Se adelanta una frase al Alcance. \\
A01-8 & \sevB & Portada & La fecha usa \verb|\today| y cambia cada vez que se compila. ``TDR'' en la portada frente a ``TdR'' en el cuerpo. & Fecha fija en la macro \verb|\fechaentrega| (a confirmar por el autor); se usa ``TdR''. \\
A01-9 & \sevB & Resumen & Comentario \texttt{\% REVISAR} sobre la cifra 40/102/160, repetida con casi la misma redacción en Resumen, Datos, Resultados y Conclusiones. & La cifra completa queda en Resumen y Resultados; Datos y Conclusiones remiten con una frase corta. \\
A01-10 & \sevB & Resumen & El uso de \texttt{run.sh} se detalla en el cuerpo del informe; por instrucción del autor pasa a un anexo separado. & El Resumen lo anuncia en una frase. \\
A01-11 & -- & Alcance & Verificado sin cambios: la cita textual de la actividad~(a) coincide palabra por palabra con el TdR; el plazo de 25 días calendario y el nombre oficial del servicio son correctos. & Se conserva. \\
\end{hallazgos}

\section{Decisiones de redacción}

\begin{itemize}[leftmargin=1.2em]
\item Las menciones al TdR se mantienen breves y en tono descriptivo (``en cumplimiento de la actividad~(a)''), sin contrastar lo exigido con lo realizado. \texttt{ssp370} se presenta como escenario intermedio adicional.
\item No se individualiza ningún modelo excluido; los excluidos se describen en conjunto (44 y 18).
\item Los modelos nombrados (FGOALS-f3-L, GISS-E2-1-G) son casos únicos de un mecanismo de corrección, no ejemplos entre muchos excluidos.
\end{itemize}

\section{Extensión}

Original: Resumen 455 palabras y Alcance 220 (675 en total). Propuesto: 698 
palabras. La reducción viene de eliminar la oración de 170 palabras sobre los bugs, que se reemplaza por una enumeración compacta; las secciones de destino conservan el detalle.

\section{Texto propuesto}

\begin{propuesto}
% ============================================================
\section{Resumen}
\label{sec:resumen}
% ============================================================

Este informe documenta el Entregable~1 del servicio de asistencia en temas de investigación del SENAMHI orientado a evaluar la influencia del cambio climático sobre El Niño-Oscilación del Sur (ENOS) en las regiones oceanográficas Niño~3.4 y Niño~1+2. En cumplimiento de la actividad~(a) de los términos de referencia, el entregable cubre la revisión de metodologías de tiempo de emergencia (TOE, por sus siglas en inglés), la recopilación de la temperatura superficial del mar (TSM) observada y la recopilación de las simulaciones históricas y futuras de los modelos CMIP6. No produce todavía resultados de desempeño, proyección ni TOE, que corresponden a los Entregables~2 a~4, pero deja construido, verificado y versionado el insumo del que esos entregables dependen.

El producto operativo es un \emph{pipeline} reproducible (repositorio \archivo{smn_toe}, con un único punto de entrada, \texttt{run.sh}) que descarga, homogeneiza, enmascara y valida la variable \texttt{tos} (TSM, tabla CMOR \texttt{Omon}) de los modelos CMIP6 que publican los experimentos \texttt{historical}, \texttt{ssp245} y \texttt{ssp585}, más \texttt{ssp370} como escenario intermedio adicional, junto con la referencia observacional ERSSTv5 de NOAA. El universo de partida son los 102 modelos CMIP6 que publican \texttt{tos}/\texttt{Omon} en la federación ESGF; 41 de ellos publican los cuatro experimentos y 40 completaron la descarga y aprobaron el control de calidad (160 combinaciones modelo-experimento, Sección~\ref{sec:resultados}). Los 62 restantes quedaron excluidos por no publicar alguno de los experimentos en ninguna fuente consultada (44) o por disponibilidad o descarga incompleta (18).

El \emph{pipeline} incorpora controles específicos frente a problemas que contaminarían el dato sin producir errores visibles. Cuatro provienen de los datos y servicios de origen: valores de relleno que el regrillado no reconoce como faltantes (FGOALS-f3-L), filtrados por valor físico; archivos que declaran grados Celsius pero contienen valores en Kelvin (GISS-E2-1-G), detectados por el valor numérico; búsquedas paginadas en ESGF que un corte de red puede truncar sin aviso, que se repiten completas; y archivos de dos variantes de grilla en un mismo modelo, que se filtran por grilla. Además, el inventario exige explícitamente los cuatro experimentos y el control de calidad se recalcula en cada corrida. El detalle de cada control está en el documento de anexos.

En paralelo, el informe desarrolla la revisión científica: la partición de incertidumbre de \citet{hawkins2012}, su adaptación por \citet{szabo2016}, la aplicación a CMIP6 y Sudamérica de \citet{bruno2023}, el precedente institucional de \citet{alvarez2024}, la caracterización del Niño costero de \citet{takahashi2019} y un estado del arte de otras metodologías de TOE. De esa revisión se deriva la elección de los dos métodos de TOE que se aplicarán en el Entregable~4: M1 \citep{szabo2016} y \emph{Time of Detection} (ToD, \citealp{gopika2024}). Las instrucciones de instalación y uso del \emph{pipeline}, así como la descripción detallada de cada paso, se entregan en un documento de anexos separado.

% ============================================================
\section{Alcance}
\label{sec:alcance}
% ============================================================

Los términos de referencia del \emph{Servicio de asistencia en temas de investigación para evaluar la influencia del cambio climático en los eventos de El Niño y sus impactos en el Perú} (en adelante, TdR) definen el primer entregable, con plazo máximo de 25 días calendario, como el informe de la actividad~(a) del numeral~6: ``Realizar la revisión de metodologías sobre tiempo de emergencia (TOE), recopilación de datos de la información de temperatura superficial del mar (TSM), y simulaciones históricas y futuras de los modelos climáticos globales CMIP6 para los escenarios SSP2-4.5 y SSP5-8.5, correspondientes a las regiones Niño~3.4 y Niño~1+2''.

Para operacionalizar este mandato, el equipo técnico lo descompone en cinco tareas: (i) revisar los antecedentes científicos (ENOS y Niño costero, partición de incertidumbre, aplicaciones previas en Sudamérica y otras metodologías de TOE); (ii) consolidar el conjunto de modelos CMIP6 que publican la variable requerida; (iii) recopilar la referencia observacional ERSSTv5; (iv) definir los dominios oceanográficos y los periodos de análisis; y (v) homogeneizar unidades, calendarios, resoluciones espaciales y máscaras océano-tierra entre todas las fuentes. La tarea~(i) es de revisión bibliográfica y se desarrolla en la Sección~\ref{sec:marco}; las tareas (ii) a (v) son computacionales y se resuelven con el \emph{pipeline} descrito en las Secciones~\ref{sec:datos} a~\ref{sec:scripts}.

El control de calidad de este entregable verifica que el dato esté completo y homogéneo, no que los modelos representen bien el clima observado: esa evaluación de desempeño es materia del Entregable~2.

\end{propuesto}

\bibliographystyle{plainnat}
\bibliography{../informe/references_rev}
\end{document}
